\documentclass[11pt]{article}

\usepackage[T1]{fontenc}
\usepackage[utf8]{inputenc}

\usepackage{lipsum}             %TODO remove

\usepackage[
  letterpaper, text={175mm,230mm}, % 215.9 x 279.4 mm   US Letter paper
  %vratio=1:1, hratio=1:1,
  text={153.33mm,230mm},         % 2:3
  %text={148mm, 222mm},           % 2:3
]{geometry}

\usepackage[spacing,tracking,letterspace=60,]{microtype}

\usepackage[dvipsnames,svgnames]{xcolor}
%\usepackage{anyfontsize} % silence font size not available warnings by autoscaling fonts
\usepackage{amssymb}            % \multimap, \mathbb
\usepackage{amsmath}
%\usepackage{latexsym}
\usepackage{stmaryrd}           % \shortrightarrow
\usepackage{mathtools}          % \dblcolon, \Coloneqq, \intertext
\usepackage{mathpartir}         % \begin{mathpar}, \infer
%\usepackage{nccmath}            % fleqn (for code environment)
\usepackage{booktabs}           % \midrule
%\usepackage{pifont}             % \ding{55}
%\usepackage{placeins}           % \FloatBarrier
%\usepackage{changepage}         % \begin{adjustwidth}
%\usepackage{etoolbox}           % \newtoggle etc
%\usepackage{graphicx}           % \rotatebox
\usepackage{url, hyperref}

\usepackage[nameinlink,noabbrev]{cleveref}
\newcommand\crefrangeconjunction{--} % use en-dashes for ranges.

%\usepackage{minted}
\usepackage{listings}           % lstlisting
\lstset{language=prolog, columns=fullflexible,
  basicstyle=\ttfamily,
  commentstyle=\color{Green},
}


%% ===== STYLING =====
\frenchspacing
\raggedbottom                   % makes things look better in the worst case

\makeatletter
\newlength\headeroffset
\setlength\headeroffset{1em}
\newcommand\makehno[1]{%
  \llap{\parbox[t]{100in}{\hfill\tbfigures#1}%
    \hspace*{\headeroffset}}%
}
\makeatother

\usepackage{titlesec}
\titleformat{\section}[block]{\scshape}{\makehno{\thesection}}{0pt}{\MakeLowercase}
%\titleformat{\section}[block]{\itshape}{\thesection\quad}{0pt}{}
%\titlespacing*{\section}{0pt}{1.5\baselineskip}{1\baselineskip}
%\titleformat{\paragraph}[runin]{\itshape}{}{}{}
\titlespacing*{\section}{0pt}{1ex}{0pt}
\titlespacing*{\paragraph}{0pt}{0pt}{1em}

\setlength\parindent{0pt}
\setlength\parskip{2ex}

%% \setlength\tabcolsep{5pt}
%% \setlength\arraycolsep{5pt}
%\def\arraystretch{1.05}


%% ===== COMMANDS =====
\newcommand\strong[1]{{\bfseries#1}}

\DeclareRobustCommand{\ensuretext}[1]{\ifmmode{\text{#1}}\else{#1}\fi}

\newcommand\todocolor{\color{IndianRed}}
\newcommand\todo[1]{\ensuretext{\todocolor #1}}
\newcommand\TODO[1]{\todo{\strong{TODO}: #1}}
\newcommand\xxx{\text{\todo{\strong{XXX}}}}
\newcommand\draft[1]{{\color{RoyalBlue}#1}}

\newcommand\fslang{\ensuremath{\lambda_{\mathrm{FS}}}}
\newcommand\catname[1]{\ensuremath{\mathrm{#1}}}
\newcommand\Set{\ensuremath{\catname{Set}}}
\newcommand\Point{\ensuremath{\catname{Set}_{*}}}
\DeclareMathOperator{\support}{supp}

\newcommand\den[1]{\llbracket{#1}\rrbracket}
\newcommand\bnfeq{\Coloneqq}
\newcommand\bnfor{\mathrel{\,|\,}}
\newcommand\defeq{\coloneqq}

\newcommand\name[1]{\ensuremath{\mathit{#1}}}
%\newcommand\<\;                 %function application space
\newcommand\<\:                 %function application space
\newcommand\fnspace{\<}         %space after a binder
\newcommand\fn[1]{\lambda{#1}.\fnspace}
\newcommand\fa[1]{\forall{#1}.\fnspace}
\newcommand\ex[1]{\exists{#1}.\fnspace}

%\newcommand\cto{\Rightarrow}    % "cartesian \to"
\newcommand\cto\to
\newcommand\lto\multimap        % "linear \to"
\newcommand\fto\Rightarrow
%\renewcommand\fto{\underset{\textsf{fin}}{\to}}           % finite map
%\newcommand\fto{\rightarrow}    % "finitely supported \to"
\newcommand\x\times
\newcommand\with{\mathbin{\&}}
\newcommand\ox\otimes
\newcommand\typename\textit
\newcommand\bool{\typename{bool}}
\newcommand\maybe{\typename{maybe}}
\newcommand\maybeof[1]{\ensuremath{\maybe~{#1}}}

%% need \usepackage{accent} for these
%\newcommand\undertype[2]{\underaccent{\textsf{\upshape\scshape#1}}{#2}}
%\newcommand\eqt[1]{\undertype{eq}{#1}}

\newcommand\G{\mathrm{\Gamma}}
\newcommand\D{\mathrm{\Delta}}
\renewcommand\O{\mathrm{\Omega}}
\renewcommand\L{\mathrm{\Lambda}}
\renewcommand\G\Gamma
\renewcommand\D\Delta
\renewcommand\O\Omega
\renewcommand\L{\Lambda}

%\newcommand\emptycx\varepsilon
\newcommand\emptycx{{\cdot}}
\newcommand\of{\mathbin{:}}
\newcommand\cxsep{\mathbin{/}}
\newcommand\hyp[2]{{#1} \of {#2}}
%\renewcommand\hyp[2]{{#1}\colon\!{#2}} % good in newtxmath[charter]
%\renewcommand\hyp[2]{{#1}\!:{#2}} %good in Euler, too loose in newtxmath[xcharter]
%% "parameter": type tagged with a variable name for use in modes
\newcommand\param[2]{\hyp{#1}{#2}}
\renewcommand\param[2]{{#1}\colon\!{#2}}

\newcommand\entails{\,\vdash\,}
%\newcommand\cJ[3]{{#3} \entails {#1} : {#2}}
%\newcommand\J[5]{{#3} \mathrel{\cdot} {#4} \cxsep {#5} \entails {#1} : {#2}}
\newcommand\Jold[5]{{#4} \cxsep {#5} \cxsep {#3} \entails {#1} : {#2}}
\newcommand\J[5]{{#3} \cxsep {#4} \cxsep {#5} \entails {#1} : {#2}}

\newcommand\rulename\textsc
\newcommand\kword[1]{\ensuremath{\textrm{#1}}}
\newcommand\const\name
%\newcommand\cname\text          %constructor name
%\newcommand\cname\name
\newcommand\cname\textrm
\newcommand\cnil{\cname{none}}
\newcommand\cjust{\cname{just}}
%\newcommand\ctrue{\cname{true}}
%\newcommand\cfalse{\cname{false}}

\newcommand\nil{\name{nil}}
\newcommand\true{\name{true}}
\newcommand\false{\name{false}}

\newcommand\letin[1]{\kword{let}~{#1}~\kword{in}~}
\newcommand\leteq[2]{\letin{{#1} = {#2}}}
\newcommand\tjust[1]{\cjust\<{#1}}
\newcommand\letjust[3]{\leteq{\tjust {#1}}{#2}{#3}}
\newcommand\letpair[4]{\leteq{({#1}, {#2})}{#3}{#4}}
\newcommand\tabsurd{\kword{absurd}\<}
\newcommand\tand{\mathrel{\kword{and}}}
\newcommand\tor{\mathrel{\kword{or}}}
%% \renewcommand\tand{~\kword{and}~}
%% \renewcommand\tor{~\kword{or}~}
%\newcommand\twhen[2]{\kword{when} \<({#1}) \<{#2}}
%\newcommand\twhen[2]{\kword{if}~ {#1} ~\kword{then}~ {#2}}
%\newcommand\twhen[2]{{#1} \Rightarrow {#2}}
%\newcommand\twhen[2]{{#2} \mathrel{\kword{when}} {#1}}
\newcommand\whenkword{\kword{when}}
\newcommand\when{\mathrel{\whenkword}}
%\newcommand\teq[2]{\name{eq} \<{#1} \<{#2}}
\newcommand\teq[2]{{#1} = {#2}}
%\newcommand\teq[2]{{#1} \mathrel{\kword{is}} {#2}}
%\newcommand\teq[2]{{#2} \gets {#1}}
\newcommand\withpair[1]{\ensuremath{\langle{#1}\rangle}}
%\newcommand\tif[3]{\kword{if}~{#1}~\kword{then}~{#2}~\kword{else}~{#3}}
\newcommand\ifjust[4]{%
  \kword{case}~{#1}\ \kword{of}\
  \cjust~{#2} \shortrightarrow {#3};\, \cnil \shortrightarrow {#4}%
}

%% \renewcommand\tand{\mathrel{\mathrm{\&\&}}}
%% \renewcommand\tor{\mathrel{\mathrm{||}}}
%% \newcommand\tand{\mathrel{\mathbf{and}}}
%% \newcommand\tor{\mathrel{\mathbf{or}}}
%% \renewcommand\ifjust[3]{%
%%   \kword{if}~\cjust~{#2} \gets {#1} ~\kword{then}~ {#3} ~\kword{else}~
%% }

\newcommand\pwild{\text{\textunderscore}}

\newcommand\typeforget{}
\newcommand\termforget[1]{#1}

\newcommand\nat{\mathbb{N}}
\newcommand\natz{\ensuremath{\nat_0}}

%% mode commands
%\newcommand\many{\top} \newcommand\mnone{0} \newcommand\mand{,} \newcommand\mplus{+}
\newcommand\many{\top}
\newcommand\mnone{\bot}
\newcommand\mand{\sqcap} \newcommand\bigmand\bigsqcap
\newcommand\mplus{\sqcup}
\newcommand\mdep[2]{{#1} \to {#2}} %\to? \shortrightarrow?
\newcommand\mundep[1]{#1} % \mundep{x} means x is ground, no dependencies
\newcommand\mex[2]{\ex{#1} {#2}}

\newcommand\mto[3]{{#1}_{#2} \to {#3}}
%\renewcommand\mto[3]{{#1}\,\{{#2}\} \to {#3}}
%\renewcommand\mto[3]{{#1} \cdot {#2} \to {#3}}
%\renewcommand\mto[3]{{#1} \star {#2} \to {#3}}
%\renewcommand\mto[3]{{#1} \oast {#2} \to {#3}}


%% %% ===== FONT OPTION 1: CHARTER ====

%% %% \usepackage[scaled=0.965, scaled=0.98, proportional,semibold]{sourcesanspro}
%% %% \usepackage[scaled=1.021,var0,varl,]{inconsolata}
%% \usepackage[scaled=0.97, osf]{XCharter}
%% \linespread{1.0333}
%% \usepackage[xcharter, alty, scaled=1.0556, % scaling makes \Rightarrow too bold :(
%%   %scaled=1.0424,
%% ]{newtxmath}


%% ===== FONT OPTION 2: PALATINO =====
\usepackage[scaled=0.96, largesc, nohelv, osf]{newpxtext}
%\usepackage[scaled=1.0176]{biolinum}
%\usepackage[scaled=0.964203422764601]{inconsolata}
\makeatletter
%\edef\zeu@Scale{.96}            %too small
\edef\zeu@Scale{.975}           %eyeballed 2026-11-08
%% HACK: newpxtext doesn't respect fontaxes' \tbfigures! :( but it _does_
%% define, eg., \tlfstyle! hm...
\usepackage{fontaxes}
\renewcommand\tbfigures{}
\renewcommand\lnfigures{}
\linespread{1.08}
\usepackage{eulervm}
%\PassOptionsToPackage{scaled=\zeu@Scale,bbscaled=\zeu@Scale,scrscaled=\zeu@Scale}{mathalfa}
\makeatother

%\renewcommand\name[1]{\ensuremath{\mathrm{#1}}}


\begin{document}

\section{notation}

\[
\begin{array}{rccll}
  \text{vars} & x,y
  \\
  \text{modes} & m,n &\bnfeq&
  \mdep{x_1 \dots\, x_n}{y}
  \bnfor \mex x m
  \bnfor m \mand n \bnfor m \mplus n
  \bnfor \many \bnfor \mnone
  \\
  \text{contexts} & \G,\L &\bnfeq& \emptycx \bnfor \G, \hyp x A
  \\
  \text{pointed contexts} & \D
  &\bnfeq& \emptycx \bnfor \D, \hyp x P
  \\
  \text{types} & A,B &\bnfeq&
  \typeforget P \bnfor A \to B \bnfor A \times B \bnfor 1 \bnfor \dots\\
  \text{pointed types} & P,Q &\bnfeq&
  P \with Q \bnfor P \ox Q \bnfor P \lto Q
  %\bnfor A \fto Q
  \bnfor \mto \G m P
  \\
  &&&%\bnfor
  \maybeof A
  %\bnfor \bool
  \bnfor \natz
  \\
  \text{expressions} & e &\bnfeq&
  \dots
  %% \termforget t
  %% \bnfor x \bnfor \fn x e \bnfor e_1 \<e_2
  %% \bnfor (e_1, e_2) \bnfor \pi_i \<e \bnfor ()
  %% \\&&&
  %% %\bnfor
  %% \ifjust{e_1}{x}{e_2}{e_3}
  %% \bnfor ...
  \\
  \text{pointed terms} & t,u &\bnfeq&
  \nil
  \bnfor x \bnfor \fn{x} t \bnfor t\<u
  \bnfor t\<x \bnfor t\<e
  \bnfor \withpair{t, u} \bnfor \pi_i \<t
  \\
  &&& (t,u) \bnfor \letpair x y t u
  \bnfor \tjust e \bnfor \letjust x t u
\end{array}
\]

\begin{tabular}{llp{86mm}}%{rp{25mm}p{110mm}}
  \it Syntax & \it Name & \it Interpretation\\
  $A,B$ & Types & Sets\\
  $A \to B$ & Functions & Functions\\
  $P,Q$ & Pointed types & Sets with a designated `zero' element, written $\nil_P$ or just $\nil$.
  For instance, booleans with $\nil = \text{false}$, or the integers with $\nil = 0$.
  \\
  $P \lto Q$ & Point-preserving maps & Functions that preserve $\nil$, that is, $f(\nil_P) = \nil_Q$.
  \\
  $A \fto P$ & Finite maps &
  Functions $f : A \to P$ such that $\support{f} \defeq \{ x : f(x) \ne \nil \}$ is finite.
\end{tabular}

\begin{tabular}{lp{108mm}}
  \em Mode & \em Interpretation \\
  $\mdep{x_1\dots\,x_n}{y}$
  &
  A finiteness dependency:
  for fixed $x_1\dots\,x_n$, finitely many $y$ satisfy the relation.
  Operationally, grounds $y$ if $x_1\dots\,x_n$ are ground.
  \\
  $\mex x m$ & Existential projection: the relation $m$ with the variable $x$ hidden.
  \\
  $m \mand n$ & Intersection of relations; union of their dependencies.
  \\
  $m \mplus n$ & Union of relations.
  Can be approximated as the intersection of the closure of their dependencies.
  \\
  $\mnone$ & Empty relation. Identity of ${\mand}$.\\
  $\many$ & Unknown; any relation. Identity of ${\mplus}$.
\end{tabular}

\begin{tabular}{lp{108mm}}
  \em Judgment & \em Interpretation
  \\
  %% $\cJ e A \G$ & $\G \to A$
  %% \\
  $\J t P \G \D m$ & Term $t$ has type $P$ using variables in $\G \cup \D$.
  It has mode $m$.
  It preserves $\nil$ in $\D$:
  for all $\hyp x Q \in \D$, we have $x = \nil_Q$ implies $t = \nil_P$.
\end{tabular}


\newcommand\sub[2]{{#1}\{#2\}}
\let\oldmodels\models
\renewcommand\models\vDash      %maybe?

\newcommand\Jmode[3]{{#1} \cxsep {#2} \entails {#3}~\text{mode}}

\pagebreak
\section{Typing rules}

\newcommand\ground[1]{\name{ground}(#1)}

Let $\ground{\G} \defeq \bigmand_{\hyp x A \in \G} \mundep{x}$ be the mode that says all variables in $\L$ are ground.

\TODO{need rule like \rulename{$\to$\,e~var} but that mixes vars and terms in argument list.}
%
\begin{mathpar}
  \infer[$\to$\,i]{
    \J t P {\G, \L} \D m
    \\
    \L = (\hyp{x_i}{A_i})_i
    \\
    \Jmode \L \emptycx n
    \\
    m \mand {\todocolor\ground{\G}} \models n
    %% \\
    %% \todo{WRONG: should be $\ground{\G}$ right?}
  }{
    \J {\fn {x_1 \dots\, x_n} t} {\mto \L n P} \G \D {\mex{x_1\dots\,x_n} m}
  }

  \infer[$\to$\,e~var]{
    \J t {\mto \L n P} \G \D m
    \\
    \L = (\hyp{x_i}{A_i})_i \subseteq \G
    \\
    \todo{need to make $n$ depend on vars used in $t$!}
    %% \fa{i} \hyp {x_i} {A_i} \in \G
  }{
    \J {t \< (x_1 \dots\, x_n)} P \G \D {m \mand n \mand \xxx}
  }
\\
  \infer[$\ox$\,i]{
    \J t P \G {\D_1} m\\
    \J u P \G {\D_2} n
  }{
    \J {(t,u)} {P \ox Q} \G {\D_1\cup\D_2} {m \mand n}
  }

  \infer[$\with$\,i]{
    \J t P \G \D m\\
    \J u Q \G \D n
  }{
    \J {\withpair{t, u}} {P \with Q} \G \D {m \mplus n}
  }

  \TODO{$\ox$\,e}

  \TODO{$\with$\,e}
\end{mathpar}

\strong{Question}: Is $n$ in \rulename{$\to$\,i} effectively $\ex{\G} m$?
\strong{Answer}: No!
We are not projecting away the variables in $\G$, we are fixing them, which is stronger.
Relational $\lambda$ is effectively a group-by, and each group receives concrete fixed values for the group-by variables in $\G$.
A concrete example:
Let $\G = (x:A)$, $\L = (y:B)$, $m = \mdep x y$.
Then $n=y$ is acceptable because $m \mand \ground{\G} = \mdep x y \mand x \models y = n$.
But $\ex{\G} m = \ex{x} \mdep x y = \many$ does not entail $y$.

\TODO{\rulename{$\to$\,i} isn't strong enough!
  Let $\G = a,b,x,y$ and $\L = x,y$.
  What if $m = \mdep x y$?
  Then what $n$ is valid?
  In fact, it's okay to let $n = \mundep y$!
  I think the rule isn't strong enough!
}

\newcommand\wantcolor{\color{RoyalBlue}}

\TODO{Does the mode judgment $\Jmode \G \D m$ actually need the $\D$ context, so that modes can refer to pointed vars in the LHSes of dependencies? Argument for no: in the type $\mto \G m P$, the mode $m$ can refer only to vars in $\G$, which cannot be pointed.

  Argument for yes: if $f : A \fto P$ is a pointed var, then $f\<x$ grounds $x$ only when $f$ is fixed, so it depends on $f$. This could come up if I compare $\fn f \fn x f\< x : (A \fto P) \lto (A \fto P)$ (obviously ok) with $\fn x \fn f f \< x : A \fto (A \lto P) \lto P$ --- this is WRONG! So what mode should $f\<x$ have?}
I've colored the mode I believe is semantically correct {\wantcolor in blue}.
%
\begin{mathpar}
  \infer{
    \J f P {\hyp x A} {\hyp f {\mto {(\param x A)} {\mundep x} P}} \many
  }{
    \J {f\<x}
       {P}
       {\hyp x A}
       {\hyp f {\mto {(\param x A)} {\mundep x} P}}
       {\many \mand \wantcolor \mdep f x}
  }
\end{mathpar}

But also, what about applying a curried function $f : \mto {(\param x A)} {\mundep x} {\mto {(\param y B)} {\mundep y} P}$?
%
\begin{mathpar}
  \infer{
    \infer{
      \J f
         {\mto {(\param x A)} {\mundep x} {\mto {(\param y B)} {\mundep y} P}}
         {\hyp x A, \hyp y B}
         {\hyp f {\mto {(\param x A)} {\mundep x} {\mto {(\param y B)} {\mundep y} P}}}
         \many
    }{
      \J {f\<x}
         {\mto {(\param y B)} {\mundep y} P}
         {\hyp x A, \hyp y B}
         {\hyp f {\mto {(\param x A)} {\mundep x} {\mto {(\param y B)} {\mundep y} P}}}
         {\xxx}
    }
  }{
    \J {f\<x\<y}
       {P}
       {\hyp x A, \hyp y B}
       {\hyp f {\mto {(\param x A)} {\mundep x} {\mto {(\param y B)} {\mundep y} P}}}
       {\xxx}
  }
\end{mathpar}

Or a curried function which doesn't ground its first argument $\hyp x A$?
%
\begin{mathpar}
  \infer{
    \infer{
      \J f
         {\mto {(\param x A)} {\many} {\mto {(\param y B)} {\mundep y} P}}
         {\hyp x A, \hyp y B}
         {\hyp f {\mto {(\param x A)} {\many} {\mto {(\param y B)} {\mundep y} P}}}
         \many
    }{
      \J {f\<x}
         {\mto {(\param y B)} {\mundep y} P}
         {\hyp x A, \hyp y B}
         {\hyp f {\mto {(\param x A)} {\many} {\mto {(\param y B)} {\mundep y} P}}}
         {\wantcolor \many}
    }
  }{
    \J {f\<x\<y}
       {P}
       {\hyp x A, \hyp y B}
       {\hyp f {\mto {(\param x A)} {\many} {\mto {(\param y B)} {\mundep y} P}}}
       {\wantcolor \mdep{f,x} y}
  }
\end{mathpar}

\section{Lipsum}

\lipsum[1-5]

\end{document}
